\chapter{Race, Coloniality, and Gender as Doxonomic Systems}
\label{ch:race-coloniality-gender}

\epigraph{The most potent weapon of the oppressor is the mind of the oppressed.}{Steve Biko}

\noindent
Race, coloniality, and gender are among the most consequential doxonomic systems in human history.
Each involves the large-scale production and institutional stabilization of beliefs about what kinds of people exist, what they are capable of, and what they deserve.
A field that claims to study how power shapes belief cannot treat these as optional examples.
They are central cases---arguably the cases where doxonomic production has been most systematic, most durable, and most materially consequential.

\section{Race as Manufactured Social Ontology}
\label{sec:ch24b-race}

\begin{definition}[Racial Regime]
\label{def:racial-regime}
\index{racial regime}
A \emph{racial regime} is a doxonomic system that produces and maintains a categorical division of humanity into hierarchically ordered groups based on perceived phenotypic or ancestral characteristics, where:
\begin{enumerate}[nosep]
  \item the categories are presented as natural rather than constructed,
  \item differential treatment is legitimated by attributed group capacities,
  \item the institutional infrastructure (law, education, media, science) continuously reproduces the classification.
\end{enumerate}
\end{definition}

The history of race is a history of doxonomic production.
Scientific racism in the 18th--20th centuries was not a marginal pseudoscience; it was a heavily funded, institutionally embedded research program that produced the ``common sense'' of racial hierarchy \autocite{dubois1903}.
Colonial administrations invested enormously in racial classification systems---census categories, legal codes, educational tracks---that made race appear as a natural fact rather than an administrative construction \autocite{said1978, fanon1961}.

\begin{model}[Racial Belief Production]
\label{mod:racial-production}
\index{racial regime!production model}
The strength of racial-hierarchy belief $\belief{R}$ is sustained by:
\[
  \frac{d\belief{R}}{dt} = \alpha_{\text{inst}} \phi_{\text{inst}} + \alpha_{\text{seg}} S(t) + \beta \belief{R}(1-\belief{R}) - \delta \belief{R} - \gamma E(t)
\]
where $\phi_{\text{inst}}$ is institutional investment in racial classification (curricula, media, law), $S(t)$ is the degree of residential and economic segregation (which generates ``evidence'' for racial difference through differential outcomes), $\beta$ is peer reinforcement, $\delta$ is natural decay, and $E(t)$ represents counter-evidence (interracial cooperation, anti-racist education, disconfirming experience).
\end{model}

The segregation term $S(t)$ is crucial: spatial and economic separation produces the differential outcomes that appear to confirm the racial narrative, creating a self-fulfilling loop analogous to the selfishness trap of Chapter~\ref{ch:human-nature}.

\section{Colonial Civilizational Narratives}
\label{sec:ch24b-coloniality}

Colonialism was, among other things, a massive doxonomic operation: the production of beliefs about civilizational hierarchy that justified extraction, governance, and violence.
\textcite{said1978} showed how ``Orientalism''---a systematic body of knowledge about the East produced by Western institutions---functioned as epistemic infrastructure for colonial rule.

The belief production chain for colonial narratives ran through:
\begin{itemize}[nosep]
  \item missionary education (replacing indigenous knowledge systems),
  \item colonial administration (legal and census categories imposing European classifications),
  \item academic disciplines (anthropology, geography, linguistics producing ``scientific'' accounts of hierarchy),
  \item media and literature (representing colonized peoples as requiring guidance),
  \item economic institutions (structuring labor markets around racial categories).
\end{itemize}

The persistence of these narratives after formal decolonization---what decolonial theorists call \emph{coloniality}---is a doxonomic phenomenon: the institutional infrastructure outlasts the political regime that built it.

\section{Gender as Doxonomic Construction}
\label{sec:ch24b-gender}

\begin{definition}[Gender Regime]
\label{def:gender-regime}
\index{gender regime}
A \emph{gender regime} is a doxonomic system producing and maintaining beliefs about:
\begin{enumerate}[nosep]
  \item natural differences between men and women in capacity, temperament, and social role,
  \item the complementarity or hierarchy of these roles,
  \item the unnaturalness of deviations from prescribed gender behavior.
\end{enumerate}
\end{definition}

Gender regimes are sustained through family socialization, religious institutions, educational tracking, media representation, workplace norms, and legal codes.
Their doxonomic cost is distributed across many institutions, making them difficult to account for with the funding-tracing methods of Chapter~\ref{ch:belief-flow-accounting}.
Much gender-regime maintenance happens through unpaid labor: mothers teaching daughters, communities enforcing norms, cultures producing representations.

This highlights a limitation of the money-centric framing: some of the most durable doxonomic systems are reproduced through institutional routine, social pressure, and embodied practice rather than through direct financial investment.
The belief production chain for gender passes heavily through layers 3--5 of the layered ontology (Principle~\ref{prin:layered-ontology}): practical schemas, identity commitments, and affective orientations, not just propositional beliefs.

\section{Intersections}
\label{sec:ch24b-intersections}

Race, gender, and class regimes do not operate independently.
They intersect, reinforce, and sometimes contradict each other.
A full doxonomic analysis must attend to how:
\begin{itemize}[nosep]
  \item racial narratives are gendered (stereotypes differ for men and women of the same racialized group),
  \item class narratives are racialized (the ``deserving poor'' is often implicitly white),
  \item gender narratives are class-inflected (the ``ideal worker'' presupposes a wife at home).
\end{itemize}

In sheaf-theoretic terms (Chapter~\ref{ch:contradiction-coherence}), the local belief systems for race, gender, and class may have different gluing obstructions depending on the institutional context.

\section{Notes and References}
\label{sec:ch24b-notes}

On race as a social construction with material effects, see \textcite{dubois1903} and \textcite{fanon1961}.
Colonial knowledge production is analyzed in \textcite{said1978}.
On the doxonomic function of education in reproducing class and racial hierarchy, see \textcite{bowles1976} and \textcite{bourdieu1984}.
For epistemic injustice and its relation to social identity, see \textcite{fricker2007} and \textcite{medina2013}.
The concept of coloniality as persisting beyond formal colonialism draws on Quijano (2000) and Mignolo (2011).
On gender as a performed and institutionally sustained category, see Butler (1990) and Connell (1987).

\section{Exercises}

\begin{exercise}
Construct a doxonomic account for racial hierarchy in a specific national context.
Identify at least four institutional channels and estimate their relative contribution to $\belief{R}$.
\end{exercise}

\begin{exercise}
Compare the doxonomic production of racial categories in two colonial contexts (e.g., British India and French West Africa).
What institutional differences led to different racial classification systems?
\end{exercise}

\begin{exercise}
Using the model of racial belief production, argue that residential segregation functions as a self-reinforcing doxonomic feedback loop.
Under what conditions can the loop be broken?
\end{exercise}

\begin{exercise}
Analyze the doxonomic infrastructure maintaining gender complementarity beliefs in a religious institutional context of your choice.
How does this infrastructure differ from the funding-driven models of earlier chapters?
\end{exercise}

\begin{exercise}
Apply the sheaf framework of Chapter~\ref{ch:contradiction-coherence} to the intersection of racial and meritocratic beliefs.
Where do the local belief systems cohere, and where do gluing obstructions arise?
\end{exercise}
